% SSSUP-FONTE-STORICA: fonti-locali/originali/2014-2020.zip
% Anno 2020, esercizio 1. Testo storico; figura ricostruita da PDF.
% Verificare editorialmente la trasposizione prima di considerarla definitiva.
\renewcommand{\Anno}{2020}
\renewcommand{\NumeroProblema}{1}

\renewcommand{\Testo}{
Un albero del Sant'Anna è cresciuto secondo le seguenti regole:
\begin{itemize}
\item il tronco è un segmento verticale di lunghezza $1$ metro ed il tronco si considera un ramo;
\item alla fine di ogni ramo di lunghezza $\ell$, dipartono due rami di lunghezza $\rho\times\ell$, con $\rho\in[0,1]$;
\item i due rami che dipartono da uno stesso ramo formano un angolo $\alpha\in[0,2\pi]$ fra di loro, ogni ramo forma quindi un angolo $\pi-\frac{\alpha}{2}$ con il ramo da cui è stato originato;
\item si definisce \emph{profondità} di un ramo il numero di giunzioni che è necessario attraversare per giungere fino al tronco. Il tronco ha quindi profondità $0$.
\end{itemize}

\begin{center}
% SSSUP-FIGURA-RICOSTRUITA: albero, ispirato a 2020figura.pdf.
% Rami disegnati ricorsivamente per alpha=pi/2 e rho=1/2.
\psset{unit=1.07cm}
\begin{pspicture}(-2.45,-.25)(2.45,4.95)
  \psset{linewidth=.65pt,linecap=1}
  \psline(0,0)(0,2.4)
  \psline(0.0000,2.4000)(-0.8485,3.2485)
  \psline(-0.8485,3.2485)(-1.4485,3.2485)
  \psline(-1.4485,3.2485)(-1.6607,3.0364)
  \psline(-1.6607,3.0364)(-1.6607,2.8864)
  \psline(-1.6607,2.8864)(-1.6076,2.8334)
  \psline(-1.6607,2.8864)(-1.7137,2.8334)
  \psline(-1.6607,3.0364)(-1.8107,3.0364)
  \psline(-1.8107,3.0364)(-1.8637,2.9834)
  \psline(-1.8107,3.0364)(-1.8637,3.0894)
  \psline(-1.4485,3.2485)(-1.6607,3.4607)
  \psline(-1.6607,3.4607)(-1.8107,3.4607)
  \psline(-1.8107,3.4607)(-1.8637,3.4076)
  \psline(-1.8107,3.4607)(-1.8637,3.5137)
  \psline(-1.6607,3.4607)(-1.6607,3.6107)
  \psline(-1.6607,3.6107)(-1.7137,3.6637)
  \psline(-1.6607,3.6107)(-1.6076,3.6637)
  \psline(-0.8485,3.2485)(-0.8485,3.8485)
  \psline(-0.8485,3.8485)(-1.0607,4.0607)
  \psline(-1.0607,4.0607)(-1.2107,4.0607)
  \psline(-1.2107,4.0607)(-1.2637,4.0076)
  \psline(-1.2107,4.0607)(-1.2637,4.1137)
  \psline(-1.0607,4.0607)(-1.0607,4.2107)
  \psline(-1.0607,4.2107)(-1.1137,4.2637)
  \psline(-1.0607,4.2107)(-1.0076,4.2637)
  \psline(-0.8485,3.8485)(-0.6364,4.0607)
  \psline(-0.6364,4.0607)(-0.6364,4.2107)
  \psline(-0.6364,4.2107)(-0.6894,4.2637)
  \psline(-0.6364,4.2107)(-0.5834,4.2637)
  \psline(-0.6364,4.0607)(-0.4864,4.0607)
  \psline(-0.4864,4.0607)(-0.4334,4.1137)
  \psline(-0.4864,4.0607)(-0.4334,4.0076)
  \psline(0.0000,2.4000)(0.8485,3.2485)
  \psline(0.8485,3.2485)(0.8485,3.8485)
  \psline(0.8485,3.8485)(0.6364,4.0607)
  \psline(0.6364,4.0607)(0.4864,4.0607)
  \psline(0.4864,4.0607)(0.4334,4.0076)
  \psline(0.4864,4.0607)(0.4334,4.1137)
  \psline(0.6364,4.0607)(0.6364,4.2107)
  \psline(0.6364,4.2107)(0.5834,4.2637)
  \psline(0.6364,4.2107)(0.6894,4.2637)
  \psline(0.8485,3.8485)(1.0607,4.0607)
  \psline(1.0607,4.0607)(1.0607,4.2107)
  \psline(1.0607,4.2107)(1.0076,4.2637)
  \psline(1.0607,4.2107)(1.1137,4.2637)
  \psline(1.0607,4.0607)(1.2107,4.0607)
  \psline(1.2107,4.0607)(1.2637,4.1137)
  \psline(1.2107,4.0607)(1.2637,4.0076)
  \psline(0.8485,3.2485)(1.4485,3.2485)
  \psline(1.4485,3.2485)(1.6607,3.4607)
  \psline(1.6607,3.4607)(1.6607,3.6107)
  \psline(1.6607,3.6107)(1.6076,3.6637)
  \psline(1.6607,3.6107)(1.7137,3.6637)
  \psline(1.6607,3.4607)(1.8107,3.4607)
  \psline(1.8107,3.4607)(1.8637,3.5137)
  \psline(1.8107,3.4607)(1.8637,3.4076)
  \psline(1.4485,3.2485)(1.6607,3.0364)
  \psline(1.6607,3.0364)(1.8107,3.0364)
  \psline(1.8107,3.0364)(1.8637,3.0894)
  \psline(1.8107,3.0364)(1.8637,2.9834)
  \psline(1.6607,3.0364)(1.6607,2.8864)
  \psline(1.6607,2.8864)(1.7137,2.8334)
  \psline(1.6607,2.8864)(1.6076,2.8334)
  \psarc[linewidth=.45pt](0,2.4){.54}{45}{135}
  \rput(0,3.12){$\alpha$}
  \rput(-.99,3.05){$\rho$}
  \rput(.99,3.05){$\rho$}
  \rput[r](-.12,1.18){$1$}
\end{pspicture}
\end{center}

\noindent\emph{Un esempio di albero con rami di profondità fino a $5$ e con $\alpha=\frac\pi2$, $\rho=\frac12$.}

\begin{enumerate}
\item[(a)] Con $\alpha=\pi$, determinare tutti i valori di $\rho$ per cui l'albero ``tocca terra'' in un punto diverso dal tronco.
\item[(b)] Con $\alpha=\frac\pi 2$, quale è l'altezza massima, in funzione di $\rho$?
\item[(c)] Dato un punto arbitrario del piano, trovare, se esiste, una coppia di valori $(\alpha, \rho)$ tale che il punto appartenga all'albero costruito con tali valori.
\end{enumerate}
}

\renewcommand{\Risposta}{
% SSSUP-BOZZA-AI: Risposta; autore=ChatGPT; data=2026-11-12; stato=DA-VERIFICARE
% SSSUP-PROVENIENZA: ChatGPT (OpenAI), sessione ProblemiSSSUP del 2026-11-12; derivato dalla traccia canonica; revisione umana pendente.
La risposta dipende dal significato di albero infinito. Ponendo il tronco
fra $(0,0)$ e $(0,1)$ e includendo i punti limite dei rami:
\begin{enumerate}
\item[(a)] $\rho\in[1/\sqrt2,1]$.
\item[(b)] $H(\rho)=(1+\rho/\sqrt2)/(1-\rho^2)$ per $0\le\rho<1$;
per $\rho=1$ l'altezza non \`e limitata.
\item[(c)] Ogni punto $(x,y)$ del piano appartiene a un albero. Una scelta \`e
\[
 M=\max\{1,|x|,|y-1|\},\qquad
 \alpha=\pi,\qquad \rho=\sqrt{\frac{M}{M+1}}.
\]
\end{enumerate}
Se si considerano soltanto i segmenti di profondit\`a finita, in (a)
l'estremo $1/\sqrt2$ va escluso e in (b), per $0<\rho<1$, $H$ \`e un
supremo non raggiunto. Anche per l'unione dei rami finiti ogni punto del
piano appartiene a qualche albero: la soluzione fornisce una costruzione
separata, che in generale usa un diverso valore di $\rho$.
}
\renewcommand{\Aiutino}{
% SSSUP-BOZZA-AI: Aiutino; autore=ChatGPT; data=2026-11-12; stato=DA-VERIFICARE
% SSSUP-PROVENIENZA: ChatGPT (OpenAI), sessione ProblemiSSSUP del 2026-11-12; derivato dalla traccia canonica; revisione umana pendente.
Prima di tutto, distinguere l'unione dei rami dalla sua chiusura.
Con $\alpha=\pi$ i rami sono alternativamente orizzontali e verticali:
separare quindi le due coordinate e sommare serie geometriche.
Con $\alpha=\pi/2$, quali sono i maggiori incrementi verticali possibili
alle profondit\`a pari e dispari? Per (c), osservare che una serie
geometrica con segni variabili pu\`o riempire un intervallo.
}
\renewcommand{\SoluzioneUfficiale}{}
\renewcommand{\SoluzioneRedazionale}{
% SSSUP-BOZZA-AI: SoluzioneRedazionale; autore=ChatGPT; data=2026-11-12; stato=DA-VERIFICARE
% SSSUP-PROVENIENZA: ChatGPT (OpenAI), sessione ProblemiSSSUP del 2026-11-12; derivato dalla traccia canonica; revisione umana pendente.
\textbf{Precisazione del modello.}
Fissiamo il tronco da $O=(0,0)$ a $(0,1)$ e il terreno sulla retta $y=0$.
Interpretiamo la ramificazione come due rotazioni di $\pm\alpha/2$
rispetto alla direzione del ramo precedente, coerentemente con la figura.
Non tronchiamo i rami quando si intersecano o attraversano il terreno.
Indichiamo con $T_{\alpha,\rho}$ l'unione di tutti i segmenti di
profondit\`a finita e con $\overline T_{\alpha,\rho}$ la sua chiusura.
Il Testo non specifica quale dei due insiemi chiami ``albero''.
La distinzione cambia alcune risposte: la manteniamo esplicita.

\textbf{(a) Contatto con il terreno.}
Per $\alpha=\pi$ i rami di profondit\`a dispari sono orizzontali,
quelli di profondit\`a pari sono verticali. Per $0\le\rho<1$, dopo
il tronco, la somma delle lunghezze dei rami verticali lungo un percorso \`e
\[
 B=\rho^2+\rho^4+\rho^6+\cdots=\frac{\rho^2}{1-\rho^2}.
\]
Ogni punto successivo al tronco ha dunque ordinata almeno $1-B$.
Se $\rho<1/\sqrt2$, questa quantit\`a \`e positiva: nessun ramo,
nemmeno al limite, tocca il terreno fuori dal tronco.

Si pu\`o scegliere un percorso che vada sempre a destra nei tratti
orizzontali e sempre verso il basso in quelli verticali. Se
$1/\sqrt2<\rho<1$, una somma finita $\rho^2+\cdots+\rho^{2k}$
supera $1$: un segmento verticale incontra quindi il terreno con
ascissa positiva. Per $\rho=1$ basta il ramo di profondit\`a $2$,
che va da $(1,1)$ a $(1,0)$.

Quando $\rho=1/\sqrt2$, ogni somma verticale finita \`e minore di $1$,
ma il percorso appena descritto converge a
\[
 \left(\frac{\rho}{1-\rho^2},\,
       1-\frac{\rho^2}{1-\rho^2}\right)=(\sqrt2,0).
\]
Questo punto appartiene alla chiusura, non all'unione dei segmenti.
Pertanto la condizione \`e $\rho\ge1/\sqrt2$ per l'albero chiuso e
$\rho>1/\sqrt2$ per l'unione dei rami finiti.

\textbf{(b) Altezza.}
Ora ogni svolta \`e di $\pi/4$. Alle profondit\`a pari la direzione
forma con la verticale un multiplo di $\pi/2$, quindi l'incremento
verticale \`e al pi\`u $\rho^{2j}$. Alle profondit\`a dispari forma
un multiplo dispari di $\pi/4$, quindi l'incremento \`e al pi\`u
$\rho^{2j+1}/\sqrt2$. Anche i punti interni ai segmenti rispettano
il limite ottenuto sommando questi incrementi positivi.
Per $0\le\rho<1$ segue
\[
 H\le\sum_{j=0}^{\infty}\rho^{2j}
      +\frac1{\sqrt2}\sum_{j=0}^{\infty}\rho^{2j+1}
   =\frac{1+\rho/\sqrt2}{1-\rho^2}.
\]
Alternando una svolta a destra e una a sinistra si ottengono
successivamente un ramo diagonale verso l'alto e un ramo verticale
verso l'alto. Tutti gli incrementi raggiungono cos\`i il rispettivo
limite: le ordinate dei vertici tendono al valore indicato.
Per $0<\rho<1$ nessuna somma finita raggiunge la somma infinita;
perci\`o si tratta di un supremo su $T_{\pi/2,\rho}$ e di un massimo
su $\overline T_{\pi/2,\rho}$. Per $\rho=0$ il massimo \`e $1$,
mentre per $\rho=1$ lo stesso percorso sale senza limite.

\textbf{(c) Un albero chiuso che contiene un punto assegnato.}
Dimostriamo prima un fatto sulle serie con segni. Se $1/2\le r<1$, allora
\[
 \left\{\sum_{j=0}^{\infty}\varepsilon_j r^j:
              \varepsilon_j\in\{-1,1\}\right\}
 =\left[-\frac1{1-r},\frac1{1-r}\right].
\]
Posto $I=[-1/(1-r),1/(1-r)]$, i due intervalli $-1+rI$ e $1+rI$
ricoprono $I$: i loro estremi esterni sono quelli di $I$ e i loro
estremi interni si sovrappongono proprio quando $r\ge1/2$.
Dato $z\in I$, si pu\`o dunque scegliere $\varepsilon_0$ con
$z=\varepsilon_0+rz_1$ e $z_1\in I$, e ripetere la scelta.
Dopo $m$ passi il resto ha modulo al pi\`u $r^m/(1-r)$, che tende a zero.
Questo dimostra il fatto e fornisce un algoritmo per scegliere i segni.

Torniamo ad $\alpha=\pi$, con $1/\sqrt2\le\rho<1$, e poniamo $r=\rho^2$.
Dopo il tronco, a ogni tratto orizzontale si pu\`o scegliere liberamente
destra o sinistra, e a ogni tratto verticale alto o basso. Le scelte
sono indipendenti in senso geometrico: entrambe le direzioni
perpendicolari sono disponibili a ogni bivio. Le coordinate dei punti
limite dei percorsi sono quindi
\[
 x=\rho\sum_{j=0}^{\infty}\varepsilon_j r^j,
 \qquad
 y=1+r\sum_{j=0}^{\infty}\eta_j r^j,
 \qquad \varepsilon_j,\eta_j\in\{-1,1\}.
\]
Applicando il fatto precedente a entrambe le coordinate, si ottiene
l'intero rettangolo
\[
 R_\rho=
 \left[-\frac{\rho}{1-\rho^2},\frac{\rho}{1-\rho^2}\right]
 \times
 \left[1-\frac{\rho^2}{1-\rho^2},1+\frac{\rho^2}{1-\rho^2}\right].
\]
Tutti i rami finiti sono contenuti nello stesso rettangolo, compreso
il tronco perch\'e $\rho^2/(1-\rho^2)\ge1$.
Pertanto $\overline T_{\pi,\rho}=R_\rho$.

Dato ora $(x,y)\in\mathbb R^2$, scegliamo
\[
 M=\max\{1,|x|,|y-1|\},\qquad \rho=\sqrt{\frac{M}{M+1}}.
\]
Allora $1/\sqrt2\le\rho<1$ e
\[
 \frac{\rho^2}{1-\rho^2}=M,
 \qquad
 \frac{\rho}{1-\rho^2}=\sqrt{M(M+1)}\ge M.
\]
Il punto assegnato appartiene dunque a $R_\rho$, che \`e l'albero
chiuso di parametri $(\pi,\rho)$.

\textbf{(c), complemento: appartenenza a un ramo finito.}
Anche senza aggiungere i punti limite, ogni punto del piano appartiene
a qualche albero. La formula precedente non basta a provarlo;
forniamo una costruzione distinta, sempre con $\alpha=\pi$.

Per $\rho=1$, l'unione dei rami contiene l'intera griglia delle rette
$x\in\mathbb Z$ e $y\in\mathbb Z$. Infatti, dopo $k$ coppie di passi
orizzontale--verticale si raggiunge ogni nodo intero $(u,v)$ con
$u\equiv v-1\equiv k\pmod2$ e $|u|,|v-1|\le k$.
Facendo crescere $k$ si raggiungono tutti i nodi con $u+v$ dispari;
da essi partono entrambi i segmenti orizzontali unitari, e dai loro
estremi entrambi i segmenti verticali unitari. Si ottengono quindi
tutti i lati della griglia. Se almeno una coordinata del punto
assegnato \`e intera, la scelta $\rho=1$ risolve il problema.

Negli altri casi possiamo supporre $x>0$, grazie alla simmetria
rispetto all'asse verticale. Scegliamo un intero pari $Y_0$ con
$|y-Y_0|<1$, e poniamo $b=Y_0-1$ e $q=|b|$. Il numero $q$ \`e dispari.
Per un intero $k$ grande, consideriamo un percorso con $2k$ tratti
orizzontali, i primi $k$ verso destra e i successivi $k$ verso sinistra.
La loro somma algebrica \`e
\[
 X_k(\rho)=\rho\,\frac{(1-\rho^{2k})^2}{1-\rho^2},
 \qquad 0<\rho<1,
\]
con estensione continua $X_k(1)=0$. Posto $\rho_0=1-x/k^2$, per
$k$ sufficientemente grande si ha $0<\rho_0<1$ e
\[
 X_k(\rho_0)\longrightarrow 2x>x.
\]
Per verificare il limite basta usare
$1-\rho_0^2\sim2x/k^2$ e $1-\rho_0^{2k}\sim2x/k$.
Per il teorema dei valori intermedi esiste dunque
\[
 \rho_k\in(1-x/k^2,1)\quad\hbox{con}\quad X_k(\rho_k)=x.
\]

Restano da scegliere i versi dei tratti verticali. Fra i primi
$2k-1$ scegliamo i primi $q$ con segno $\operatorname{sgn}(b)$,
e quelli rimasti a coppie con segni $+1,-1$.
Il numero dei tratti rimasti \`e pari. Per $\rho=1$ l'ordinata
raggiunta sarebbe esattamente $1+b=Y_0$.
Per $\rho=\rho_k$, posto $r=\rho_k^2$, l'ordinata $Y_k$ all'inizio
dell'ultimo tratto verticale soddisfa
\[
 |Y_k-Y_0|\le(1-r)\left(\frac{q(q+1)}2+k\right)\longrightarrow0.
\]
Infatti $1-r^j\le j(1-r)$ controlla i primi $q$ termini;
ciascuna coppia successiva vale $r^j-r^{j+1}$ e ha modulo
al pi\`u $1-r$. Inoltre $1-r<2x/k^2$.
L'ultimo tratto verticale ha lunghezza $\rho_k^{4k}\to1$.
Poich\'e $|y-Y_0|<1$, per $k$ abbastanza grande risulta
\[
 |y-Y_k|<\rho_k^{4k}.
\]
Scegliendo l'ultimo tratto verso l'alto o verso il basso, secondo
il segno di $y-Y_k$, esso contiene il punto $(x,y)$.
Per $x<0$ si riflettono tutte le scelte orizzontali.

La costruzione \`e finita: si aumenta $k$, si sceglie una radice di
$X_k(\rho)=|x|$ nell'intervallo indicato e si verifica l'ultima
disuguaglianza. Le stime dimostrano che per $k$ sufficientemente grande
il procedimento riesce. Il punto appartiene allora a un ramo di
profondit\`a $4k$, non soltanto alla chiusura dell'albero.
Non viene cercato il valore minimo della profondit\`a.
}
\renewcommand{\SoluzioniAlternative}{}
\renewcommand{\Osservazioni}{}
\renewcommand{\Commenti}{}
